% !TEX root = ../main.tex
%-----------------------------------------------------------------------
%  Motivation.
%  \input by ../main.tex -- not compiled on its own.
%-----------------------------------------------------------------------
%-----------------------------------------------------------------------
\section{Motivation}
\label{sec:motivation}
%-----------------------------------------------------------------------

The adjoint returns the gradient of a scalar cost functional with respect to
every control in one reverse sweep. That is what makes it the right tool, and
also what makes it expensive. In the DINO configuration as it stands:

\begin{itemize}[nosep]
  \item a five-year adjoint sweep on 27 MPI ranks takes $23\,\mathrm{h}$
        wall-clock, about $622$ core-hours, and writes $8.3\,\mathrm{GB}$;
  \item the gradient holds for \emph{one} background state, \emph{one}
        parameter setting and \emph{one} cost functional;
  \item the cost-functional section indices are compiled in as Fortran
        \code{parameter} statements in \code{cost\_atlantic\_heat.F}, so moving
        the section means a rebuild rather than a namelist edit.
\end{itemize}

So any question about how the sensitivity pattern itself changes with the
mixing parameters costs one full adjoint run per setting. That is the
bottleneck I want to remove. If the surrogate works, it buys us four things:

\begin{enumerate}[nosep]
  \item parametric sensitivity and uncertainty quantification across the
        mixing-parameter space, at something like $\mathcal{O}(10^{-6})$ of the
        present cost per additional sample;
  \item the ability to re-target the cost functional --- a different latitude, a
        different depth range --- without recompiling and re-running;
  \item gradients inside a calibration or optimisation loop, where calling the
        true adjoint at every iteration is out of the question, with the true
        adjoint kept for periodic correction;
  \item some interpretability, in that a surrogate which generalises across
        parameter settings has learned something about which features of the
        background state control the sensitivity pattern.
\end{enumerate}

The underlying question is how the sensitivity of the basin's overturning
transport to the mixing parameterisation depends on the ocean state and on the
parameters themselves.

